\documentclass{article}
\usepackage{decomp}
\begin{document}
\sourcefile{elife\_figure\_supplements/fig1-figsupp1.jpg … fig2-figsupp5.jpg and elife\_api\_article.json}
\decompnote{The eLife PDF of this article prints only the titles of the eight figure supplements (main.tex, after the Figure 1 and Figure 2 legends). The images here are the full-size JPEGs from the eLife IIIF image server, and the legends (title and caption) are taken verbatim from the eLife API record of the article (\texttt{elife\_api\_article.json}, objects \texttt{fig1s1} … \texttt{fig2s5}, fields \texttt{label}, \texttt{title}, \texttt{caption}); see the folder's \texttt{SOURCE.txt}. HTML markup was converted to LaTeX (bold panel letters, italics, sub- and superscripts); links were reduced to their text. There are no page numbers. The check against the JSON is \texttt{work/check\_supp\_legends.py} (README).}

\decompnote{Figure 1—figure supplements 1–3.}

\begin{center}
\includegraphics[width=0.6\linewidth,height=0.75\textheight,keepaspectratio]{figures/sfig1s1.png}
\end{center}
\noindent\textbf{Figure 1—figure supplement 1.} As reported previously (Mikula and Denk, 2015, Figure 4b), X-ray imaging exhibits a contrast similar to scanning electron microscopy (SEM). In the following μCT reslices and X-ray projection images of \emph{en bloc} stained brain tissue are compared to the corresponding serial-section SEM image from the same tissue block. (\textbf{a-g}) Covariation of the standardized μCT and standardized SEM intensity in the same tissue area. The μCT and SEM pixel intensities were averaged along the vertical direction (black arrow) in the regions of interest (ROI) shown in e+f/k+l. (\textbf{b-h}) μCT reslice and (\textbf{c-i}) SEM overview image of the same tissue section. The black rectangle indicates the location of the ROI shown in e+f/k+l. (\textbf{d-j}) Scatter plot of the average SEM pixel intensity versus the average μCT reslice pixel intensity in the ROI shown in e+f/k+l, averaged along the vertical direction (black arrow) (d: \emph{r}=0.923, p<1e-9, j: 0.944, p<1e-9; Pearson correlation coefficient). (\textbf{m}) Covariation of the standardized X-ray and standardized SEM intensity as a function of cortical depth. The X-ray intensity was measured along the red line in the X-ray projection image shown in (n). The SEM intensity was measured along the blue line in the SEM image shown in (o) after averaging the intensity values along the horizontal axis (indicated by the blue arrow). The intensities were standardized by subtracting the mean and dividing by the standard deviation. Note, the X-ray intensity corresponds to a projection through the entire 4 mm brain punch, whereas the SEM intensity has been only accumulated over the width of about 1.5 mm of the original punch in b, because the embedded block was trimmed to a width of about 1.5 mm for ultra-thin section collection. (\textbf{n}) X-ray projection image of a 4 mm punch of brain tissue that was incubated in 2\% buffered OsO\textsubscript{4} for 20 hours. (\textbf{o}) Low-resolution SEM image of a 40 nm thick section (same as in Figure 1—figure supplement 3) from the sample in n, after dehydration with ethanol and acetonitrile and embedding in hard Epon. The section was collected orthogonally to the transverse plane of the punch at a depth of about 1 mm from the sample surface. The image has been rescaled to account for the shrinkage of the sample during the dehydration and embedding as well as for compression from sectioning. (\textbf{p}) Scatter plot of the average SEM pixel intensity versus the X-ray projection image pixel intensity (\emph{r}=0.243, p<4.45e-9; Pearson correlation coefficient). Scale bar 100 μm.
\decompnote{Panel ranges ``a-g'', ``b-h'', ``c-i'', ``d-j'' are printed so (they mean a and g, b and h, and so on). In (o), ``same as in Figure 1—figure supplement 3'' probably means Figure 1—figure supplement 2, which shows the SEM images of that section; supplement 3 has only plots.}
\alttext{Sixteen panels in three blocks (a--f, g--l, m--p), each block arranged as a plot (left) with images (right). (a) Line plot, ``Standardized (a.u.)'' (−3.5 to 1.5; ticks −3 to 1) against ``Distance (μm)'' (0 to about 1320; ticks 0, 400, 800, 1200), boxed axes; ``SEM'' (blue, noisy) and ``μCT'' (orange, smooth) both rise from about −2.7 (SEM) and −3.3 (μCT) at 0 μm to about 0 by 200 μm, reach about 1 near 700--850 μm and fall to about −2 near 1300 μm. (b) Greyscale μCT reslice of the tissue block (dark tissue, light background) with a long black rectangle labelled ``e''; small scale bar at lower right. (c) The matching SEM overview image (light tissue) with a black rectangle labelled ``f''. (d) Scatter plot of ``mean SEM intensity (a.u.)'' (about 186--223; ticks 190--220) against ``mean μCT intensity (a.u.)'' (about 44--103; ticks 50--100): small black dots forming two roughly parallel rising bands from about (45, 191) to (102, 220). (e), (f) Long horizontal strips: the ROIs of the μCT reslice and the SEM image, each with a black downward arrow at the right end and a black scale bar at lower right. (g) As (a) for a second ROI: ``Standardized (a.u.)'' (−3 to 3) against ``Distance (μm)'' (0--1200); SEM and μCT rise together from about −2.6 (μCT) and −1.8 (SEM) to about 1.5 (μCT) and 2.2 (SEM) at about 1070 μm, SEM much noisier. (h), (i) μCT reslice and SEM overview of another section with a horizontal black rectangle labelled ``k'' and ``l'', and a light region (ventricle) at lower left; scale bars at lower right. (j) Scatter plot, ``mean SEM intensity (a.u.)'' (about 145--170) against ``mean μCT intensity (a.u.)'' (100--180): a dense rising cloud of black dots from about (112, 147) to (177, 169). (k), (l) The ROI strips, with downward arrows and scale bars. (m) Line plot, ``Standardized (a.u.)'' (−5 to 2) against ``Depth (μm)'' (0--1600): SEM (blue) starts at about −0.6, peaks at about 1.4 near 300 μm, falls to about −1.2 by 1150 μm and ends at about −1.5; μCT (orange) rises steeply from about −4.4 to about 0.7 by 450 μm, has a bump to about 0.85 near 850 μm and is flat near 0 after 1200 μm. (n) Greyscale X-ray projection image of the punch, dark domed top, with a short red line running down from the top surface. (o) SEM image of the section with a blue vertical line along the left edge and a blue arrow pointing right. (p) Scatter plot, ``mean SEM intensity (a.u.)'' (about 55--120) against ``mean X-ray intensity (a.u.)'' (50--100): dots form a curve rising slowly from about (50, 72) to (85, 100), then steeply to about (90, 112), and a near-vertical branch falling from about (92, 110) to (88, 57).}
\figuredata{a}{figures/sfig1s1a.csv}{figures/sfig1s1.py}{digitized from the eLife JPEG (tier 3) by sfig1s1\_extract.py}
\figuredata{d}{figures/sfig1s1d.csv}{figures/sfig1s1.py}{digitized from the eLife JPEG (tier 3): every dark pixel of the dot cloud}
\figuredata{g}{figures/sfig1s1g.csv}{figures/sfig1s1.py}{digitized from the eLife JPEG (tier 3)}
\figuredata{j}{figures/sfig1s1j.csv}{figures/sfig1s1.py}{digitized from the eLife JPEG (tier 3): every dark pixel of the dot cloud}
\figuredata{m}{figures/sfig1s1m.csv}{figures/sfig1s1.py}{digitized from the eLife JPEG (tier 3)}
\figuredata{p}{figures/sfig1s1p.csv}{figures/sfig1s1.py}{digitized from the eLife JPEG (tier 3): every dark pixel of the dot cloud}

\begin{center}
\includegraphics[width=0.95\linewidth,height=0.55\textheight,keepaspectratio]{figures/sfig1s2.png}
\end{center}
\noindent\textbf{Figure 1—figure supplement 2.} Ultrastructural preservation at various tissue depths after 20 hr of X-ray-assisted staining with 2\% buffered OsO\textsubscript{4}. First, a 4 mm punch of aldehyde-fixed brain tissue was immersed in 2\% buffered OsO\textsubscript{4} monitored using X-ray assisted staining for 20 hr. Afterwards, the sample was washed, dehydrated with ethanol and acetonitrile and embedded in hard Epon. Next, 40-nm-thick sections of the trimmed sample block were collected orthogonally to the transverse plane of the brain punch at a depth of about 1 mm from the tissue surface. Electron microscopy images of the section were acquired at low magnification (left) as well as high magnification at different cortical depths (middle +right column) using scanning electron microscopy. Scale bars: Left: 100 μm, middle column: 2 μm, right column: 200 nm.
\alttext{SEM micrographs (backscattered-electron contrast, membranes dark). Left: a tall low-magnification overview of the section from the cortical surface (top, with a light, bulging area above the tissue) downward: grey neuropil with lighter round cell bodies and white blood vessel profiles, a band of darker speckled tissue across the middle-lower part, and lighter tissue below; three small dark squares (near the top, the middle and the bottom) are joined by black lines to the middle column; scale bar at lower left. Middle column: three medium-magnification images (top, middle, bottom depth), each with a small dark-blue square marking the area enlarged at the right and a scale bar at lower left: dense neuropil with continuous dark membranes, mitochondria and small processes; the middle image has many dark, thick, irregular myelin-like profiles running horizontally; the bottom image has an elongated vessel-like profile with a thick dark wall. Right column: the three high-magnification images: (top) neurites with sharp dark membranes, vesicle-filled profiles and a mitochondrion with visible cristae; (middle) a large profile with a very thick dark ring (myelin) and a dark band across the top; (bottom) neurites with continuous membranes and a dark rounded organelle. Short scale bars at lower left.}

\begin{center}
\includegraphics[width=0.6\linewidth,height=0.75\textheight,keepaspectratio]{figures/sfig1s3.png}
\end{center}
\noindent\textbf{Figure 1—figure supplement 3.} Spatial intensity profiles of accumulated heavy metals after t=20 hr of incubation (left column) and temporal tissue expansion profiles (right column) for different osmium reduction protocols in comparison to 2\% buffered OsO\textsubscript{4} (blue, n=13) and buffered solution of 2\% OsO\textsubscript{4} reduced with 2.5\% potassium ferricyanide K\textsubscript{4}[Fe(CN)\textsubscript{6}] (green, n=6) (mean ±s.d.). All solutions were buffered by cacodylate. (\textbf{a-b}) Buffered solution of 2\% OsO\textsubscript{4} reduced with 2.5\% potassium ferricyanide K\textsubscript{3}[Fe(CN)\textsubscript{6}] (red, n=2). (\textbf{c-d}) Buffered solution of 2.5\% potassium ferrocyanide K\textsubscript{4}[Fe(CN)\textsubscript{6}] after 20 hr of incubation in buffered 2\% OsO\textsubscript{4} solution (light green, n=6). (\textbf{e-f}) Buffered solution of 2.5\% potassium ferricyanide K\textsubscript{3}[Fe(CN)\textsubscript{6}] after 20 hr of incubation in buffered 2\% OsO\textsubscript{4} solution (orange, n=4). Note, the intensity for K\textsubscript{3}[Fe(CN)\textsubscript{6}] in (\textbf{e}) appears to be about 8\% higher than for OsO\textsubscript{4} alone. But as can be seen in (f), the tissue height shrinks by about 2\% within the 20 hr of incubation. Under the assumption that the shrinkage is isotropic in all three dimensions, this corresponds to a volume reduction of about 5.9\% which can explain the observed increase in the intensity.
\decompnote{``potassium ferricyanide K\textsubscript{4}[Fe(CN)\textsubscript{6}] (green, n=6)'' in the title is printed so; the green series is the ferrocyanide-reduced osmium (Ferro-redOs, K\textsubscript{4}[Fe(CN)\textsubscript{6}]), as the legend under the panels says.}
\alttext{Six plots in three rows, each row with a legend below it. Left column (a, c, e): ``Intensity (a.u.)'' against ``Depth (μm)'' (about −150 to 1200; ticks 0, 500, 1000); y labels 0.4 and 1 only, a thin horizontal line at 1 and an unlabelled tick at the top of the axis (about 1.27). Right column (b, d, f): ``Change in sample height (\%)'' (−10 to 20) against ``Time (min)'' (0 to about 1210), thin line at 0. In every panel the same two reference series are drawn: ``OsO\textsubscript{4}'' (blue; intensity about 0.9 from 600 μm on; height +5.5 \% at 1200 min) and ``2.5\% (Ferro) K\textsubscript{4}[Fe(CN)\textsubscript{6}] + 2\% OsO\textsubscript{4}'' (dark green; intensity peak about 1.13 near 600 μm with a very wide band reaching above the axis; height −4.2 \%). Row 1 legend adds ``2.5\% (Ferri) K\textsubscript{3}[Fe(CN)\textsubscript{6}] + 2\% OsO\textsubscript{4}'' (red): in (a) it rises more slowly, to about 0.98 at 600 μm and a peak of about 1.06 near 950 μm, then drops to about 0.87 at 1200 μm, with a wide pink band; in (b) it stays near 0 (about −0.5 to +0.7 \%). Row 2 adds ``2.5\% (Ferro) K\textsubscript{4}[Fe(CN)\textsubscript{6}] (after 2\% OsO\textsubscript{4})'' (light green): in (c) it rises slowly and almost linearly to about 0.83 at 950 μm; in (d) it stays near −0.5 \% to about 600 min and rises to about +3 \% at 1200 min. Row 3 adds ``2.5\% (Ferri) K\textsubscript{3}[Fe(CN)\textsubscript{6}] (after 2\% OsO\textsubscript{4})'' (orange): in (e) it follows the blue curve closely but slightly higher (about 0.94--0.95 from 600 to 1000 μm); in (f) it falls to about −2 \% by 500 min and stays there.}
\figuredata{a}{figures/sfig1s3a.csv}{figures/sfig1s3.py}{digitized from the eLife JPEG (tier 3) by sfig1s3\_extract.py}
\figuredata{b}{figures/sfig1s3b.csv}{figures/sfig1s3.py}{digitized from the eLife JPEG (tier 3)}
\figuredata{c}{figures/sfig1s3c.csv}{figures/sfig1s3.py}{digitized from the eLife JPEG (tier 3)}
\figuredata{d}{figures/sfig1s3d.csv}{figures/sfig1s3.py}{digitized from the eLife JPEG (tier 3)}
\figuredata{e}{figures/sfig1s3e.csv}{figures/sfig1s3.py}{digitized from the eLife JPEG (tier 3)}
\figuredata{f}{figures/sfig1s3f.csv}{figures/sfig1s3.py}{digitized from the eLife JPEG (tier 3)}

\decompnote{Figure 2—figure supplements 1–5.}

\begin{center}
\includegraphics[width=0.95\linewidth,height=0.55\textheight,keepaspectratio]{figures/sfig2s1.png}
\end{center}
\noindent\textbf{Figure 2—figure supplement 1.} Long incubation in potassium ferrocyanide dissolves and disintegrates the tissue. (\textbf{a}) Whole mouse brain from Figure 1d after long incubation in osmium. (\textbf{b}) Subsequently, the osmium solution was replaced by 2.5\% buffered potassium ferrocyanide solution and the sample was incubated for 16 hr. The tissue clearing is particularly well visible in the cerebellum and olfactory bulbs (arrowheads). (\textbf{c}) After another 108 hr in buffered potassium ferrocyanide a large fraction of the cerebellum and the olfactory bulbs has been dissolved. Scale bar 2 mm.
\decompnote{``Whole mouse brain from Figure 1d'' is printed so; the whole mouse brain is shown in Figure 1g.}
\alttext{Three greyscale X-ray projection images of a whole mouse brain in side view (olfactory bulb at the left, cerebellum at the right), each with a ``2mm'' scale bar at lower left, and a grey-scale key ``Intensity (a.u.)'' from 0 (white) to 1 (black) at lower right. (a) The brain is uniformly mid-grey to dark, darkest at the upper left (frontal region). (b) Lighter overall; a black arrowhead and a white arrowhead at the upper left point at the frontal/olfactory region, and a black arrowhead at the right points at the cerebellum, which is much lighter than in (a). (c) Lighter still; the outline at the upper left is ragged and dark-streaked (black arrowhead) and the cerebellum at the right (black arrowhead) is faint and partly missing.}

\begin{center}
\includegraphics[width=0.95\linewidth,height=0.55\textheight,keepaspectratio]{figures/sfig2s2.png}
\end{center}
\noindent\textbf{Figure 2—figure supplement 2.} Experimentally measured change in heavy metal accumulation of a sample immersed in double distilled H\textsubscript{2}O (ddH\textsubscript{2}O). Prior to washing the sample with ddH\textsubscript{2}O the sample has been stained in 2\% buffered OsO\textsubscript{4} for 22 hr. The first row shows the temporal profile of the change in heavy metal accumulation at different cortical depths, while the second row shows the spatial profile of the change in heavy metal density at different times after washing onset. The black circles indicate the location of axon tracts. The third row shows the measured spatial-temporal change of heavy metal accumulation (left) and the corresponding distribution of heavy metal accumulation in projection images of the sample at the beginning of the experiment and after about 22.5 hr of incubating in double distilled H\textsubscript{2}O (right).
\alttext{A MATLAB figure page titled ``2\%Os1.wash, IAC20220409\_B20220406'', in the same layout as the source-data pages (source\_data.tex). Top row: six plots titled ``x=200 μm'' … ``x=1200 μm'', ``Intensity (a.u.)'' (ticks 0.5--0.8) against ``Time t (min)'' (0--1200): each blue trace starts at a peak (about 0.77 at 200 μm, about 0.83--0.88 deeper) and decays within about 300 min to a slowly falling level (about 0.69 at 200 μm, 0.75 at 400 μm, 0.78--0.80 deeper at 1200 min). Middle row: six plots titled ``t=30 min'' … ``t=1200 min'', intensity against ``Depth x (μm)'': each rises steeply from about 0.5 at 0 μm to about 0.8 by 500 μm, has a small peak near 850--900 μm marked by a black circle (about 0.88 at 30 min, 0.82 at 1200 min) and dips slightly after it; the curves are slightly lower at later times. Bottom row: colour map ``data'' (Depth x vs Time t, both 0--1200) with colour bar ``Intensity (a.u.)'' 0.45--0.85: blue below about 100 μm, green to yellow-orange above, nearly constant in time; and two projection images ``t=0min'' and ``t=1338.4833min'' in false colour, almost identical.}
\figuredata{all line plots}{figures/sfig2s2\_traces.csv}{figures/sfig2s25.py}{digitized from the eLife JPEG (tier 3) by sfig2s25\_extract.py}
\figuredata{colour map}{figures/sfig2s2\_map\_data.csv}{figures/sfig2s25.py}{read through the colour bar from the eLife JPEG (tier 3)}

\begin{center}
\includegraphics[width=0.7\linewidth,height=0.4\textheight,keepaspectratio]{figures/sfig2s3.png}
\end{center}
\noindent\textbf{Figure 2—figure supplement 3.} (\textbf{a}) Spatial intensity profile of accumulated heavy metals at the beginning t=0 hr and after t=20 hr of immersion in double distilled H\textsubscript{2}O (n=1). Note, the slight decrease in intensity can be explained on one hand by the wash out of unbound osmium from the sample as well as by expansion of the tissue which reduces the local heavy metal density. (\textbf{b}) Temporal tissue expansion profile for 20 hr immersion in double distilled H\textsubscript{2}O (n=1).
\alttext{(a) ``Intensity (a.u.)'' (labels 0.4 and 1, thin line at 1, an unlabelled tick at the top of the axis) against ``Depth (μm)'' (about −150 to 1200; ticks 0, 500, 1000). Two thick curves: ``t = 0 hrs after in H\textsubscript{2}O'' (blue) and ``t = 20 hrs after in H\textsubscript{2}O'' (black); both start at about 0.44 left of 0 μm and rise steeply after 0 μm; blue reaches about 0.84 at 500 μm and 0.86 at 900 μm (small peak) and ends at about 0.83; black is lower, about 0.76 at 500 μm, 0.82 at 900 μm and 0.78 at 1200 μm. (b) ``Change in sample height (\%)'' (−10 to 20, line at 0) against ``Time (min)'' (0--1200): one black curve rising from about 0.6 \% to about 2 \% at 200 min and about 3 \% from 600 min on. Legend below (a).}
\figuredata{a}{figures/sfig2s3a.csv}{figures/sfig2s3.py}{digitized from the eLife JPEG (tier 3) by sfig2s3\_extract.py}
\figuredata{b}{figures/sfig2s3b.csv}{figures/sfig2s3.py}{digitized from the eLife JPEG (tier 3)}

\begin{center}
\includegraphics[width=0.95\linewidth,height=0.4\textheight,keepaspectratio]{figures/sfig2s4.png}
\end{center}
\noindent\textbf{Figure 2—figure supplement 4.} Effective diffusion coefficient D\textsuperscript{e}, unmasking rate constant k\textsuperscript{unmask}, binding rate constant k\textsuperscript{on}, the OsO\textsubscript{4} concentration in the solution C\textsuperscript{0}, the initial density of binding sites S\textsuperscript{0}, the initial density of masked binding sites M\textsuperscript{0} and the sample curvature height H for n=13 different 4 mm punches that have been immersed for 22 hr in 2\% buffered OsO\textsubscript{4} solution. The parameters have been fitted to the experimental data by simulating the model Equations 2–5; 7, while C\textsuperscript{0} has been measured for each sample individually.
\alttext{Bar chart with individual values, y axis 0--10 (ticks 0, 5, 10; no axis title). Seven groups, labelled below: ``D\textsuperscript{e} in 10\textsuperscript{4} μm\textsuperscript{2} min\textsuperscript{-1}'', ``k\textsuperscript{unmask} in 10\textsuperscript{-5} mm\textsuperscript{3} nmol\textsuperscript{-1} min\textsuperscript{-1}'', ``k\textsuperscript{on} in 10\textsuperscript{-4} mm\textsuperscript{3} nmol\textsuperscript{-1} min\textsuperscript{-1}'', ``C\textsuperscript{0} in 10\textsuperscript{2} nmol mm\textsuperscript{-3}'', ``S\textsuperscript{0} in 10\textsuperscript{2} nmol mm\textsuperscript{-3}'', ``M\textsuperscript{0} in 10\textsuperscript{2} nmol mm\textsuperscript{-3}'', ``H in 10\textsuperscript{2} μm''. Each group has light-grey dots (individual samples, jittered sideways) left of a black bar with an upward error bar. Read off the plot: D\textsuperscript{e} bar about 3.4, error bar to about 5.5, dots from about 1.6 to 4.8 and one at about 9.6; k\textsuperscript{unmask} about 1.9 (to 2.7), dots 1.2--2.6 and one at 3.9; k\textsuperscript{on} about 3.3 (to 4.3), dots 1.8--4.6; C\textsuperscript{0} about 0.78 (to 0.85), dots in a tight row at about 0.8; S\textsuperscript{0} about 3.5 (to 3.85), dots 2.9--4.0; M\textsuperscript{0} about 1.4 (to 1.8), dots 1.1--2.3; H about 3.7 (to 4.8), dots 1.8--5.2.}
\figuredata{all groups}{figures/sfig2s4.csv}{figures/sfig2s4.py}{digitized from the eLife JPEG (tier 3) by sfig2s4\_extract.py; overlapping dots merge}

\begin{center}
\includegraphics[width=0.95\linewidth,height=0.55\textheight,keepaspectratio]{figures/sfig2s5.png}
\end{center}
\noindent\textbf{Figure 2—figure supplement 5.} Experimentally measured accumulation of heavy metals (blue traces) in buffered 2\% OsO\textsubscript{4} solution of a 4 mm punch of aldehyde-fixed liver tissue (n=1). The corresponding simulated model was fitted to the experimental data of the liver sample (black traces). The first row shows the temporal profile of osmium accumulation at different depths measured from the top of the sample, while the second row shows the spatial profile of osmium accumulation at different times after staining onset. The third row shows the measured spatial-temporal accumulation of osmium (left), the corresponding fitted model simulations (middle) and the distribution of osmium in projection images of the sample at the beginning of the experiment and after about 22 hr of incubation (right). The fitted model parameters were: D\textsuperscript{e} = 1.496·10\textsuperscript{5} μm\textsuperscript{2}/min, k\textsuperscript{unmask} = 1.681·10\textsuperscript{–5}mm\textsuperscript{3}/nmol/min, k\textsuperscript{on} = 2.190·10\textsuperscript{–4}mm\textsuperscript{3}/nmol/min, S\textsuperscript{0}=2.997·10\textsuperscript{2} nmol/mm\textsuperscript{3}, M\textsuperscript{0}=1.901·10\textsuperscript{2} nmol/mm\textsuperscript{3}, H=8·10\textsuperscript{2} μm, C\textsuperscript{0}=0.823·10\textsuperscript{2} nmol/mm\textsuperscript{3}. The effective diffusion coefficient appears to be >4 X larger in liver tissue compared to brain tissue. The residual standard error of the fitted model was SE\textsubscript{res} <0.018.
\alttext{A MATLAB figure page titled ``OsLiver, IAC20211203\_L20211201'', in the source-data layout. Top row: six plots ``x=200 μm'' … ``x=1200 μm'', ``Intensity (a.u.)'' (ticks 0.5--0.9) against ``Time t (min)'' (0--1200): blue data and black model rise steeply in the first 100--200 min and then slowly, to about 0.74 at 200 μm, 0.82 at 400 μm, 0.87 at 600 μm and 0.91--0.93 at 800--1200 μm; the model is slightly above the data at 400--1000 μm. Middle row: six plots ``t=30 min'' … ``t=1200 min'', intensity against ``Depth x (μm)'': at 30 min a broad hump (about 0.6 at 300--500 μm) falling to 0.5 at 1200 μm; with time the profiles rise and the hump moves deeper, until at 1200 min the curve rises from about 0.5 at 0 μm to about 0.92 at 800--1200 μm; the model curves have a kink at about 800 μm. Bottom row: colour maps ``data'' and ``model: R\textsuperscript{2}\textsubscript{adj}=0.98661, SE\textsubscript{res}=0.017913'' (Depth x vs Time t, 0--1200), colour bars ``Intensity (a.u.)'' 0.5--0.9, both blue at the left edge and turning yellow-orange from the top-right; and projection images ``t=0min'' (a pale blue outline of the domed sample on a dark-blue background) and ``t=1351min'' (the sample yellow-orange inside, with a small round blue-green spot in the middle).}
\figuredata{all line plots}{figures/sfig2s5\_traces.csv}{figures/sfig2s25.py}{digitized from the eLife JPEG (tier 3) by sfig2s25\_extract.py}
\figuredata{colour maps}{figures/sfig2s5\_map\_data.csv, figures/sfig2s5\_map\_model.csv}{figures/sfig2s25.py}{read through the colour bars from the eLife JPEG (tier 3)}

\end{document}
